% GENERATED FILE. Edit canonical lessons, then run npm run generate:books.
% canonical-source-hash: fnv1a64:4ed1f84890f82ee6
% canonical-lessons: ES-C545-criar, ES-C545-ninez, ES-C545-juventud, ES-C545-adolescencia, ES-C545-generacion

\chapter{De niño}
\label{ch:es-de-nino}
\hlchaptermodality{545}

\begin{chapteropening}
\textbf{By the end of this chapter:} \emph{I can talk about growing up: who brought me up, my childhood, my youth and my generation.}

Complete the last lesson of chapter 545: I can talk about growing up: who brought me up, my childhood, my youth and my generation.
\end{chapteropening}

\section[criar]{criar — to bring up, to raise}

\label{lesson:ES-C545-criar}



\begin{quote}
Before the new one: say the Spanish for the pregnancy, then the Spanish for the newborn baby.
\end{quote}

\subsection*{You'll want to know: criar}
\textbf{criar} — \textquotedblleft{}to bring up, to raise\textquotedblright{}.

\begin{quote}
Mi abuela me crió en el pueblo.
\end{quote}

\begin{quote}
\emph{My grandmother brought me up in the village.}
\end{quote}

\emph{Criarse} is to grow up: \emph{Me crié en Valencia.} It takes an accent: \emph{crío}.

\begin{cousinweb}
\emph{criar} is inherited from Latin \emph{creāre} 'to make, to bring forth' (EtymDB). Spanish kept the everyday sense of rearing children; \emph{crear} 'to create' is the same verb taken again from books.
\end{cousinweb}

\begin{grammarlens}[title={What you've built}]
One more word for telling, explaining and describing at B1.
\end{grammarlens}

\subsection*{Your turn}
\begin{itemize}
\raggedright
  \item \emph{Say it:} \emph{criar}
  \item \emph{Say it:} \emph{criar}, once more
  \item \emph{Say it:} a sentence of your own with \emph{criar}
  \item \emph{Recall:} say the Spanish for the pregnancy, then the Spanish for the newborn baby, then say \emph{criar} again
\end{itemize}

\begin{tcolorbox}[breakable,colback=teal!4,colframe=teal!35!black,title={Before you move on}]
What does criar mean? (\textquotedblleft{}To bring up, to raise\textquotedblright{}.) Say it once more.
\end{tcolorbox}
\section[la niñez]{la niñez — childhood}

\label{lesson:ES-C545-ninez}



\begin{quote}
Before the new one: say the Spanish for the newborn baby, then the Spanish for to bring up.
\end{quote}

\subsection*{You'll want to know: la niñez}
\textbf{la niñez} — \textquotedblleft{}childhood\textquotedblright{}.

\begin{quote}
Tuve una niñez muy feliz.
\end{quote}

\begin{quote}
\emph{I had a very happy childhood.}
\end{quote}

\emph{La infancia} means the same.

\begin{cousinweb}
\emph{niñez} is built on \emph{niño} 'child' with the ending \emph{-ez}. \emph{Niño} is usually traced to a baby-talk word, \emph{ninno}, not to any Latin noun.
\end{cousinweb}

\begin{grammarlens}[title={What you've built}]
One more word for telling, explaining and describing at B1.
\end{grammarlens}

\subsection*{Your turn}
\begin{itemize}
\raggedright
  \item \emph{Say it:} \emph{la niñez}
  \item \emph{Say it:} \emph{la niñez}, once more
  \item \emph{Say it:} a sentence of your own with \emph{la niñez}
  \item \emph{Recall:} say the Spanish for the newborn baby, then the Spanish for to bring up, then say \emph{la niñez} again
\end{itemize}

\begin{tcolorbox}[breakable,colback=teal!4,colframe=teal!35!black,title={Before you move on}]
What does la niñez mean? (\textquotedblleft{}Childhood\textquotedblright{}.) Say it once more.
\end{tcolorbox}
\section[la juventud]{la juventud — youth; young people}

\label{lesson:ES-C545-juventud}



\begin{quote}
Before the new one: say the Spanish for to bring up, then the Spanish for childhood.
\end{quote}

\subsection*{You'll want to know: la juventud}
\textbf{la juventud} — \textquotedblleft{}youth; young people\textquotedblright{}.

\begin{quote}
En mi juventud viajé mucho.
\end{quote}

\begin{quote}
\emph{In my youth I travelled a lot.}
\end{quote}

\emph{La juventud} is also young people as a group.

\begin{cousinweb}
\emph{juventud} comes from Latin \emph{iuventūs}, \emph{iuventūtis} 'youth', built on \emph{iuvenis} 'young' (Lewis \& Short). English \emph{juvenile} is from the same root.
\end{cousinweb}

\begin{grammarlens}[title={What you've built}]
One more word for telling, explaining and describing at B1.
\end{grammarlens}

\subsection*{Your turn}
\begin{itemize}
\raggedright
  \item \emph{Say it:} \emph{la juventud}
  \item \emph{Say it:} \emph{la juventud}, once more
  \item \emph{Say it:} a sentence of your own with \emph{la juventud}
  \item \emph{Recall:} say the Spanish for to bring up, then the Spanish for childhood, then say \emph{la juventud} again
\end{itemize}

\begin{tcolorbox}[breakable,colback=teal!4,colframe=teal!35!black,title={Before you move on}]
What does la juventud mean? (\textquotedblleft{}Youth; young people\textquotedblright{}.) Say it once more.
\end{tcolorbox}
\section[la adolescencia]{la adolescencia — adolescence, the teenage years}

\label{lesson:ES-C545-adolescencia}



\begin{quote}
Before the new one: say the Spanish for childhood, then the Spanish for youth.
\end{quote}

\subsection*{You'll want to know: la adolescencia}
\textbf{la adolescencia} — \textquotedblleft{}adolescence, the teenage years\textquotedblright{}.

\begin{quote}
Pasé la adolescencia en Barcelona.
\end{quote}

\begin{quote}
\emph{I spent my teenage years in Barcelona.}
\end{quote}

\emph{En la adolescencia} — in your teens.

\begin{cousinweb}
\emph{adolescencia} is a learned borrowing from Latin \emph{adolēscentia} 'youth' (EtymDB), built on \emph{adolēscere} 'to grow up'.
\end{cousinweb}

\begin{grammarlens}[title={What you've built}]
One more word for telling, explaining and describing at B1.
\end{grammarlens}

\subsection*{Your turn}
\begin{itemize}
\raggedright
  \item \emph{Say it:} \emph{la adolescencia}
  \item \emph{Say it:} \emph{la adolescencia}, once more
  \item \emph{Say it:} a sentence of your own with \emph{la adolescencia}
  \item \emph{Recall:} say the Spanish for childhood, then the Spanish for youth, then say \emph{la adolescencia} again
\end{itemize}

\begin{tcolorbox}[breakable,colback=teal!4,colframe=teal!35!black,title={Before you move on}]
What does la adolescencia mean? (\textquotedblleft{}Adolescence, the teenage years\textquotedblright{}.) Say it once more.
\end{tcolorbox}
\section[la generación]{la generación — the generation}

\label{lesson:ES-C545-generacion}



\begin{quote}
Before the new one: say the Spanish for youth, then the Spanish for adolescence.
\end{quote}

\subsection*{You'll want to know: la generación}
\textbf{la generación} — \textquotedblleft{}the generation\textquotedblright{}.

\begin{quote}
Tres generaciones viven en esa casa.
\end{quote}

\begin{quote}
\emph{Three generations live in that house.}
\end{quote}

\begin{cousinweb}
\emph{generación} is a learned borrowing from Latin \emph{generātiō} 'a begetting' (EtymDB), from \emph{generāre}, built on \emph{genus} 'birth, kind'.
\end{cousinweb}

\begin{grammarlens}[title={What you've built}]
One more word for telling, explaining and describing at B1.
\end{grammarlens}

\subsection*{Your turn}
\begin{itemize}
\raggedright
  \item \emph{Say it:} \emph{la generación}
  \item \emph{Say it:} \emph{la generación}, once more
  \item \emph{Say it:} a sentence of your own with \emph{la generación}
  \item \emph{Recall:} say the Spanish for youth, then the Spanish for adolescence, then say \emph{la generación} again
\end{itemize}

\begin{tcolorbox}[breakable,colback=teal!4,colframe=teal!35!black,title={Before you move on}]
What does la generación mean? (\textquotedblleft{}The generation\textquotedblright{}.) Say it once more.
\end{tcolorbox}
